\setcounter{chapter}{1}

\chapter{Organizzazione}
    \par
        L'Università degli Studi di Udine offre i suoi insegnamenti in due sedi distaccate, situate rispettivamente a Pordenone e a Gorizia. \newline
        Entrambe fanno riferimento alla Direzione Centri e Servizi (DICS) e svolgono funzioni di supporto alle attività di Ateneo, alla gestione delle carriere degli studenti e ai servizi alla didattica, attraverso strutture dell'Amministrazione Centrale (AMCE) o dei Dipartimenti. \newline
        Le segreterie dei due Centri lavorano a stretto contatto con gli Enti Territoriali promuovendo l'Ateneo sul territorio e più in generale coltivando i rapporti con gli altri Enti pubblici e privati esistenti. \newline
        Entrambe operano secondo le indicazioni del Responsabile DICS e dei Delegati del Rettore, sulla base dei Regolamenti di Funzionamento\footnote{\cite{IstruzioniCEGO-CEPO}}. \newline
        Nonostante i compiti diversificati, ogni membro della Direzione del CEGO lavora all'unisono per garantire la continuità del servizio che viene offerto dall'Università. 
        Durante il tirocinio da me svolto ho avuto l'occasione di interfacciarmi con una realtà complessa come può esserlo l'organizzazione di un Polo Universitario, avvicinandomi alla componente umana dell'organizzazione del CEGO. 
% ---------------------------------------------------------------------------------------
    \section{Organigramma del Centro Polifunzionale di Gorizia}
        \par
            Ciascuna intervista al personale del CEGO è stata condotta utilizzando due domande di riferimento: "Qual è il suo ruolo all'intero del CEGO?" e "In relazione alla domanda precedente, quali mansioni ricopre il suo ruolo?". 

            Per maggior chiarezza, si ritiene opportuno esplicitare alcune righe: 
            \begin{itemize}
                \item \textit{DICS}: DIrezione Centro e Servizi;
                \item \textit{AMCE}: AMministrazione Centrale;
                \item \textit{DIAF}: DIrezione Amministrativa e Finanza;
                \item \textit{ADID}: Area servizi per la DIDattica;
                \item \textit{ASTU}: Area servizi agli STUdenti;
                \item \textit{ALOG}: Area edilizia e LOGistica;
                \item \textit{DISO}: ufficio DIritto allo StudiO;
                \item \textit{DAGS}: Direzione Acquisti, Gare e Servizi in outsorucing;
                \item \textit{DARU}: Direzione Risorse Umane e Affari generali;
                \item \textit{DIPC}: DIrezione Pianificazione, Controllo e valutazione;
                \item \textit{DIP}: DIPartimento; 
                \item \textit{CCS}: Consiglio del Corso di Studio;
                \item \textit{CDD}: Consiglio Di Dipartimento;
                \item \textit{CdS}: Corso di Studio;
                \item \textit{L}: Laurea Triennale; 
                \item \textit{LM}: Laurea Magistrale;
                \item \textit{SPES}: Servizio di Posta Elettronica per gli Studenti;
                \item \textit{Scheda SUA}: Scheda Unica Annuale;
            \end{itemize}
            Un elenco più completo si può reperire nella pagina web \cite{AcronimiSigle} dell'Università degli Studi di Udine. \newline
            Da precisare che il CEGO opera in due poli diversi, il primo è ubicato presso Via Santa Chiara 1/A mentre il secondo, distante pochi metri, si trova nel Corso Verdi al civico 4. Entrambi sono situati nel centro di Gorizia. 
% ---------------------------------------------------------------------------------------
        \subsubsection{Prof.ssa Renata Kodilja - Delegata del Rettore}
            La Prof.ssa Renata Kodilja, Delegata del Rettore, ricopre anche la delega alla Direzione del CEGO, che corrisponde alla presa in carico delle relazioni con il territorio, nello specifico con gli Enti e le Istituzioni del Sistema Gorizia che concorrono a  finanziare la presenza dell'Università nel territorio. Inoltre si occupa della gestione del funzionamento interno dell'aspetto di \textsl{governance} e delle strategie d'uso degli spazi. Nell'eventualità che si verifichi una carenza degli spazi, la Delegata cerca soluzioni e provvede a contratti e a convenzioni che coprano il fabbisogno dell'utenza accolta  dall'Ateneo\footnote{Questo implica il fatto che il numero massimo di immatricolazioni dipenda direttamente dalla disponibilità delle aule. Ad oggi il numero massimo di matricole corrisponde a 230, cui si aggiungono 20 matricole di provenienza extra Europea.}. \newline
            Per quanto riguarda le relazioni con il territorio questo significa invece governare i patti e gli accordi o le convenzioni con gli Enti che sostengono, finanziano, hanno interesse, motivazione e attenzione per far sì che l'Università sia a Gorizia. A queste interlocuzioni partecipano due controparti principali: il Consorzio per lo Sviluppo del Polo Universitario di Gorizia (CONSUNIGO), di cui è presidente l'Avv. Sergio Orzan\footnote{\cite{CONSUNIGO}}, e il Consorzio Universitario di Pordenone (UNIPORDENONE), di cui è presidente il Dott. Paolo Candotti. Questi Consorzi ospitano le due Sedi dell'Università degli Studi di Udine in Friuli-Venezia Giulia, mentre la Sede di Gemona non usufruisce del sostegno di alcun Consorzio. \newline
            CONSUNIGO e UNIPORDENONE raccolgono i finanziamenti e si interfacciano con la Sede Universitaria per sostenerne la didattica e il funzionamento. Il CONSUNIGO gestisce in buona misura i contratti per le docenze temporanee, mentre le spese per il mantenimento e il funzionamento delle strutture sono gestite direttamente dall'Università. \newline
            Un altro finanziatore importante è la Cassa di Risparmio di Gorizia (CARIGO) la cui missione è quella di sviluppo e promozione del territorio, mantenendo la compresenza dei due Atenei della Regione. 
            Si aggiungono la Camera di Commercio, Industria, Artigianato e Agricoltura Venezia Giulia (CCIAA Venezia Giulia Trieste Gorizia)\footnote{\cite{CCIAA}}, che accorpa le Camere di Commercio di Venezia e Trieste; i Gruppi Europei di Cooperazione Territoriale (GECT)\footnote{\cite{GECTGOEZTS}} per la città di Gorizia, ossia il consorzio tra le città di Gorizia, Nova Gorica e San Pietro-Vertoiba, che accoglie e gestisce i fondi europei per lo sviluppo e l'integrazione transfrontaliera, finanziando progetti di integrazione transfrontaliera e progetti Interreg EU.
% ---------------------------------------------------------------------------------------
        \subsubsection{Dott.ssa Cristina Prizzi - Funzionaria Amministrativa (Capocentro)}
            La Dott.ssa Cristina Prizzi, Funzionaria Amministrativa, ha due ruoli principali: supportare e collaborare con la Prof.ssa Kodilja in tutte le sue attività da Delegata del Rettore, nello specifico nei rapporti con enti pubblici, nei rapporti con i finanziatori, in presenza a determinati eventi e occupandosi della stesura di verbali. La sua figura funge anche da ponte tra gli uffici della Sede Centrale e il resto del CEGO presente nel polo di Santa Chiara. È la responsabile e la coordinatrice del personale, di conseguenza della sede, verificandone il buon andamento. Sempre in relazione agli enti, funge anche come punto di riferimento, trasmettendo le richieste di finanziamento e altro materiale. \newline
            In secondo ruolo, la Dott.ssa Prizzi svolge anche mansioni di categoria organizzativa, lavorando a più che stretto contatto con la Delegata per monitorare l'effettivo svolgimento dei vari compiti.
% ---------------------------------------------------------------------------------------
        \subsubsection{Dott.ssa Cristina Delmonaco - Referente del Corso di Laurea Relazione Pubbliche (CL-20)}
            La Dott.ssa Cristina Delmonaco si occupa della didattica e della logistica del Corso di Laurea Relazioni Pubbliche, corso tenuto solamente nella Sede di Gorizia; nello specifico organizza l'attribuzione delle aule agli insegnamenti e i calendari degli appelli d'esame. \newline
            Inoltre si occupa, in parallelo alla carica di Referente del Corso di Laurea Relazioni Pubbliche, della gestione organizzativa e logistica degli eventi che si svolgono in corso d'anno nel polo di via Santa Chiara o di Corso Verdi. \newline
            In collaborazione con la Dott.ssa Erika Bader, si occupa in parte della Segreteria Studenti. 
% ---------------------------------------------------------------------------------------
        \subsubsection{Dott. Carlo Carratù - Referente amministrativo dei Corsi di Laurea}
            Il Dott. Carlo Carratù si occupa dell'aspetto finanziario per il CEGO e per l'area DICS; gestisce gli ordinamenti didattici e, a cascata, tutto ciò che riguarda le attività normative, ad esempio il piano di studi e i relativi allegati, nonché le eventuali modifiche. Si occupa dei Consigli dei Corsi di Studio, delle Commissioni Qualità, degli orari e degli appelli d'esame per il Corso di Laurea Discipline dell'audiovisivo, dei media e dello spettacolo (DAMS; CL-3) nonché dell'inserimento nei vari moduli U-GOV e della scheda SUA. Quest'ultima è uno strumento gestionale del Ministero dell'Università e della Ricerca, funzionale alla progettazione, alla realizzazione, all'autovalutazione e alla ri-progettazione del CdS. La scheda mira a definire la domanda di formazione, esplicitare l'offerta formativa, certificare i risultati di apprendimento, chiarire i ruoli e le responsabilità che attengono alla gestione del sistema di Assicurazione della Qualità dell'Ateneo oltre che riesaminare periodicamente l'impianto del CdS e i suoi effetti per apportare le necessarie modifiche\footnote{\cite{schedaSua}}. 
% ---------------------------------------------------------------------------------------
        \subsubsection{Dott.ssa Cristina Feresin - Ufficio Tirocini}
            La Dott.ssa Cristina Feresin segue l'organizzazione dei tirocini curricolari e volontari per gli studenti, fornisce le indicazioni sulle procedure e sulle pratiche di tirocini dal primo inserimento in azienda fino alla chiusura, funge da referente per i tirocini interni e fornisce anche un \textsl{helpdesk} per i tirocini, rivolto sia agli studenti sia alle aziende. Inoltre organizza l'incontro obbligatorio per i tirocini del Corso di Laurea Relazioni Pubbliche. La Dott.ssa è anche membro della Commissione Tirocini, lavorando in stretta collaborazione con l'Ufficio Tirocini della Sede di Udine.\newline
            Infine promuove tra gli studenti gli avvisi relativi alla ricerca di personale altamente qualificato da parte delle aziende e fornisce supporto alla Segreteria della Direzione. 
% ---------------------------------------------------------------------------------------
        \subsubsection{Dott. Andrea Castelli - Supporto IT}
            Il Dott. Andrea Castelli fornisce assistenza tecnica soprattutto agli Uffici Amministrativi ma anche al resto dei dipendenti, dei professori, degli studenti. L'assistenza fornita è generale e copre ogni aspetto dell'ambito informatico d'Ateneo: reti internet, computer, stampanti, scanner, etc. Inoltre installa i nuovi software utili al corretto e migliore andamento dei dispositivi.
% ---------------------------------------------------------------------------------------
        \subsubsection{Dott. Francesco Barreca - Tecnico Multimediale}
            Il Dott. Francesco Barreca lavora come tecnico multimediale e specialista in didattica a distanza dell’Università degli Studi di Udine. Il Dott. Progetta e segue la manutenzione degli impianti multimediali delle aule, occupandosi di tutto ciò che riguarda la proiezione, i monitor, l'impianto audio e i sistemi per la didattica online; gestisce la piattaforma di e-learning dell'Ateneo basata su Moodle e la piattaforma dedicata alle registrazioni delle lezioni \emph{uniud.ubicast.tv}; collabora alla realizzazione di eventi istituzionali come l'Inaugurazione dell'Anno Accademico, la Cerimonia del Tocco, lo Student Day e il PhD Day. \newline
            Inolte si occupa della coordinazione tra le ditte esterne in outsourcing, le quali forniscono supporto tecnico per convegni e conferenze; tali ditte si occupano anche della manutenzione delle aule didattiche e curano l'installazione di nuovi impianti multimediali negli spazi di attuali o future sedi universitarie.
% ---------------------------------------------------------------------------------------
        \subsubsection{Dott. Lino Saro - Funzionario Informatico}
            Il Dott. Lino Saro si occupa della didattica laboratoriale, dell'allestimento dei laboratori informatici compresi i server CUDA, della piattaforma Moodle Learning Management System (LMS) e dei vari servizi e-learning. Prepara la correzione automatica dei vari programmi, è membro della Commissione dell'Alfabetizzazione Informatica per gli Studenti. \newline
            In parte, gestisce l'amministrazione del servizio Active Directory (AD) di Ateneo, i domini e le varie tabelle DNS; costruisce i vari bandi per i concorsi che utilizzano le piattaforme online LMS per entrare a lavorare in università: nello specifico questo si traduce nel preparare gli account temporanei per l'accesso ai test d'ammissione. Struttura sia la prova stessa sia il corso; tutta la preparazione viene anche sottoposta ad un controllo incrociato con gli account che poi verranno utilizzati per fare il test. \newline
            Lavora a stretto contatto con colleghi sviluppatori di software in base alle esigenze dell'Ateneo, utilizzando principalmente linguaggi del tipo \textsl{C\#} mentre i colleghi gestiscono applicativi del tipo Esse3 e Xform su richiesta dei Dipartimenti e dell'AMCE, verificando le firme digitali per la didattica e la gestione dei protocolli. Infine il Dott. Saro collabora con i Tecnici che gestiscono i server e i vari dispositivi delle reti e reti wireless, firewall e sicurezza. 
% ---------------------------------------------------------------------------------------
        \subsubsection{Dott.ssa Grazia Puiatti - Bibliotecaria}
            La Dott.ssa Grazia Puiatti è referente della Biblioteca della Sede di Gorizia, quindi si occupa dei prestiti e delle restituzioni dei volumi intrateneo e intrabiblioteche, entrambi tramite la Piattaforma Alma. Sovrintende alla catalogazione, all'inventario della Sede, alla gestione di ordini magazzino e di cancelleria. Tramite la piattaforma Udiscover informa l'utenza del catalogo multimediale e fornisce consulenza bibliografica. \newline
            Gestisce inoltre gli utenti interni ed esterni, preparando e posizionando il materiale informativo. Si occupa della gestione di Tutor e dei vari Collaboratori Studenteschi, lavorando a stretto contatto con la Dott.ssa Feresin per il coordinamento dei tirocinanti.
% ---------------------------------------------------------------------------------------
        \subsubsection{Gabriele Marcatti - Aiuto Bibliotecario}
            Il Sig. Gabriele Marcatti si interfaccia con gli utenti della Biblioteca di Gorizia, gestendo i prestiti e le restituzioni, sia quelli nella Sede sia quelli intrateneo che quelli intrabibliotecari, tramite la Piattaforma Alma. Controlla il corretto allestimento degli scaffali, dell'archivio, delle etichette per le collocazioni, delle timbrature dei libri e infine della stampa del materiale informativo. 
% ---------------------------------------------------------------------------------------
        \subsubsection{Dott.ssa Erika Bader - Referente Segreteria Studenti}
            La Dott.ssa Erika Bader segue la carriera degli studenti dall'orientamento in entrata e le immatricolazioni fino alla compilazione della domanda di Laurea e alla restituzione del Tesserino; aiuta nella compilazione dei Piano di Studio, rilascia i certificati di partecipazione ad eventi, interviene negli incontri mirati alla stesura di Tesi di Laurea e fornisce supporto all'accesso dei test d'ingresso.
% ---------------------------------------------------------------------------------------
        \subsubsection{Prof. Andrea Mariani - Coordinatore DAMS}
            Il Prof. Andrea Mariani coordina il Corpo Docenti e monitora l'andamento della didattica e delle carriere. Svolge anche il ruolo di Docente del Corso di Studi in "Teorie Critiche dei Media e dei Nuovi Media" ed è delegato all'orientamento e all'internazionalizzazione. \newline
            In parallelo al ruolo di Coordinatore e Docente, il Prof. è anche ricercatore in "Archeologia dei Media"\footnote{\cite{AndreaMariani}}. È responsabile e direttore dei progetti che si svolgono nel laboratorio di \textsl{Digital Storytelling} presso la Sede del Dipartimento di Studi Umanistici e del Patrimonio Culturale (DIUM). Tali progetti possono coinvolgere gli studenti come tirocinanti ma sono principalmente mirati all'attività di formazione, organizzata in workshop di vario livello, rivolta a insegnanti, a professionisti e a lavoratori che vogliano acquisire strumenti di comunicazione efficaci e innovativi\footnote{\cite{DigitalStorytellingLAB}}.
            Lavora a stretto contatto con il Dott. Carratù per la costruzione dell'orario didattico delle lezioni e della pianificazione degli appelli d'esame, oltre che per la calendarizzazione della didattica integrativa comprensiva di seminari ed eventi. Collabora con la Dott.ssa Feresin per la gestione dei tirocini curricolari. 
% ---------------------------------------------------------------------------------------
    \section{Differenze tra CEPO e CEGO}
            Le differenze principali tra CEPO e CEGO si possono dividere in due categorie principali: la gestione dei Regolamenti interni e la gestione della Didattica, ma non mancano anche alcuni tratti in comune. \newline
            Le divergenze e convergenze sono state estrapolate dal documento \cite{IstruzioniCEGO-CEPO}. 
% ---------------------------------------------------------------------------------------
        \subsubsection{Gestione dei Regolamenti interni - CEPO}
            La prima principale differenza tra CEPO e CEGO è la gestione dei regolamenti interi, in quanto il CEPO ha mantenuto il regolamento di Funzionamento e, di conseguenza, ha mantenuto l'organo del Consiglio Direttivo. L'aver mantenuto il Regolamento comporta una differenza quasi sostanziale in quanto ne definisce la funzione di organi di competenze e definisce al Consiglio Direttivo delle mansioni da delegare. Altre differenze tra i due Centri polifunzionali si possono trovare in svariati ambiti, tra cui la contabilità e la gestione ordinaria e straordinaria degli immobili. 
% ---------------------------------------------------------------------------------------
        \subsubsection{Gestione della Didattica - CEPO}
            Le gestioni Didattiche del CEPO, come del CEGO, sono in collaborazione funzionale con i Dipartimenti di riferimento dei vari CdS; tuttavia, alcuni di questi CdS sono gestiti in coordinazione con i direttori dei DIP. Va precisato che al CEPO sono presenti diversi CdS tra cui alcuni tenuti da altri Atenei. \newline
            Nello specifico i CdS dell'Università di Udine a Pordenone sono: 
            \begin{itemize} 
                \item \textsl{L - Tecnologie digitali e comunicazione per le industrie creative}: il CEPO ne gestisce il coordinamento funzionale con il Dipartimento delle scienze Matematiche, Informatiche e Fisiche (DMIF), ne gestisce l'orario degli esami e delle lauree nonché il calendario degli appelli, il CEPO ne monitora le pratiche studenti; 
                \item \textsl{LM - Comunicazione Multimediale e Tecnologie dell'Informazione}: CdS gestito in coordinato funzionale con il DMIF; 
                \item \textsl{L - Infermieristica}: CdS gestito dal Dipartimento di MEDdicina (DMED), mentre il CEPO si occupa dei tirocini, dei servizi di supporto alla didattica e alla definizione dell'offerta formativa, oltre al monitoraggio delle pratiche studenti; 
                \item \textsl{L - Ingegneria Industriale dell'Energia}: il CEPO si occupa delle attività di supporto alla didattica in sede, collaborando con il Dipartimento Politecnico di Ingegneria e Architettura (DPIA) e monitorando le pratiche studenti; 
                \item \textsl{LM - Banca e Finanza}: il CEPO si occupa delle attività di supporto alla didattica in sede, collaborando con il Dipartimento di Scienze Economiche e Statistiche (DIES);
            \end{itemize}
            Le equipollenze vengono riconosciute formalmente per tutti i CdS offerti al CEPO: la documentazione necessaria passa attraverso l'organo dell'ASTU che procede con la verifica dei Titoli di Studio. 
        \par
            Il CEPO viene sempre informato dei vari CCS e dei lavori del CDD in caso di impegni a carico del Centro stesso e di cambiamenti facenti riferimento ai vari CdS offerti in sede.\newline
            Ponendo ora il focus i CdS in Tecnologie digitali e comunicazione per le industrie creative e in Comunicazione Multimediale e Tecnologie dell'Informazione, il CEPO segue tutto l'iter delle relative pratiche studenti e procede con la stesura dei verbali di CCS. Inoltre il Centro procede all'informazione, tramite i canali ufficiali, dell'emanazione di bandi relativi alle supplenze e supporta i Docenti nelle possibili problematiche riguardo l'uso della piattaforma Esse3, oltre che nella risoluzione di ogni tipo di problema tecnico si dovesse riscontrare; nel caso di problemi riferiti agli immobili, il CEPO utilizza il \cite{PortaleSelfServiceSysAid}. Per questi due CdS il CEPO definisce anche i calendari degli appelli d'esame e gli Appelli di Laurea con le relative Commissioni, nonché la loro pubblicazione sulle pagine web, dopo essersi confrontato con il DMIF. Per quanto riguarda le Discussioni di Laurea, queste si possono verificare in due spazi principali: in centro città a Palazzo Montereale Mantica, previo consenso del Consorzio, oppure nella sede di Via Prasecco. \newline
            Una precisazione va aggiunta per la LM in Comunicazione Multimediale, in quanto il corso collabora alla gestione del \textit{Double Degree} con l'Università di Klagenfurt, occupandosi al supporto della Commissione didattica per la selezione dei candidati e la pubblicazione dei risultati, nonché il successivo riconoscimento dei Crediti Formativi Universitari conseguiti all'estero. \newline
            Il Centro di Pordenone, infine, si interfaccia con il Centro Linguistico di Ateneo per l'organizzazione di test di piazzamento in sede per le lingue straniere e, successivamente, per l'organizzazione dei corsi delle lingue straniere. Il CEPO ne supporta l'operato per eventuali prove di accertamento organizzate dai CdS. 
% ---------------------------------------------------------------------------------------
        \subsubsection{Gestione dei Regolamenti interni - CEGO}
            Come accennato in precedenza, la gestione dei Regolamenti interni è uno dei due punti fondamentali dove i due centri differiscono; il CEGO, pur essendo una sede a sé stante con un'autonomia diversa rispetto al Centro di Pordenone sotto alcuni punti di vista, fa riferimento direttamente all'Ateneo dato che non utilizza più il regolamento da Centro Polifunzionale. Questo implica una gestione interna degli spazi disponibili tramite i servizi di prenotazione dell'Ateneo e un coordinamento continuo con la Diocesi presente nel polo di Corso Verdi, essendo sotto contratto d'affitto. \newline
            Per quanto riguarda invece le discussioni delle Lauree, queste vengono discusse solamente nel polo di Santa Chiara, lasciando la procedura e la gestione degli spazi completamente in carico al CEGO. 
% ---------------------------------------------------------------------------------------
        \subsubsection{Gestione della Didattica - CEGO}
           L'offerta formativa del CEGO comprende diversi corsi gestiti in coordinazione con i DIP di riferimento, nello specifico: 
            \begin{itemize}
                \item \textsl{L - Relazioni Pubbliche}: fa riferimento al Dipartimento di Lingue e Letterature, Comunicazione, Formazione e Società (DILL);
                \item \textsl{LM - Comunicazione Integrata per le Imprese}: fa riferimento al Dipartimento di Lingue e Letterature, Comunicazione, Formazione e Società (DILL);
                \item \textsl{L - Discipline dell'audiovisivo, dei media e dello spettacolo (DAMS)}: fa riferimento al Dipartimento di studi Umanistici e del Patrimonio Culturale (DIUM).
            \end{itemize}
            I tre CdS sono gestiti dalla Segreteria Didattica in quasi piena autonomia, occupandosi della stesura dei CCS, supportando le attività formative in relazione all'offerta, inserendo i vari regolamenti didattici nella Scheda SUA-CdS e, come nel caso di Pordenone, supportando i docenti nelle eventuali problematiche della piattaforma Esse3. Inoltre il CEGO collabora con la Segreteria ADID per la redazione delle guide e dei Manifesti agli studi dei CdS, aggiustando eventualmente parti del manifesto decise da CCS e/o CDD. \newline
            Dal punto di vista delle lezioni dei vari CdS e dei rispettivi appelli d'esame, la Segreteria Didattica si occupa della creazione e dell'organizzazione dei calendari, pubblicando poi il tutto sulle varie piattaforme. La Segreteria provvede anche della stesura del calendario degli Appelli di Laurea, calendario che viene poi pubblicato sulla piattaforma Esse3. Gli spazi prenotati per lezioni, esami e Lauree vengono gestiti tramite la piattaforma Easyroom. 
% ---------------------------------------------------------------------------------------
        \subsubsection{Caratteristiche in comune e differenze minori}
            \par
                Partendo dalle caratteristiche in comune dei due Centri, entrambi i Responsabili provvedono alle richieste d'acquisto e alla consegna di beni e servizi utili al funzionamento. I Responsabili inoltre sono nominati annualmente agenti contabili del fondo amministrativo assegnato alla rispettiva struttura. 
            \par
                In caso di lavori di manutenzione, periodici o straordinari che siano, è necessario richiedere questi interventi tramite un sistema di ticket fornito dall'università\footnote{\cite{PortaleSelfServiceSysAid}}. L'interfaccia di questo portale differisce tra studente e personale interno mostrando diverse categorie disponibili. Nel caso di uno studente, il portale mostra le seguenti categorie: Esse3, Servizi per gli Studenti, Ricercatori e Dottorati, Tirocini di Medicina, Servizi Informatici e Sistema Qualità di Ateneo. Mentre, nel caso del personale interno, il portale mostra le seguenti categorie: Servizi Informatici, Esse3, DARU, Servizi Tecnici, Area Edilizia e Sistema Qualità di Ateneo. \newline 
                Per quanto riguarda gli interventi di manutenzione edilizia, ogni intervento va comunicato al proprietario della struttura del rispettivo Centro; quindi nel caso del CEPO questi interventi vanno comunicati e concordati con il Consorzio Universitario, mentre nel caso del CEGO bisogna differenziare i due poli: in via Santa Chiara, l'edificio è stato donato dal comune all'Ateneo e quindi gli interventi sono concordati con l'Ateneo stesso, mentre per Corso Verdi ogni intervento, indipendentemente dalla natura, va comunicato e concordato con la Diocesi di Gorizia, in quanto il CEGO è legato all'edificio tramite contratto d'affitto.
            \par
                Entrambi i Centri offrono il servizio di Segreteria Studenti in collaborazione con la segreteria ADID. La comunicazione tra le diverse Segreterie avviene sia in forma scritta sia in presenza in riunioni periodiche. Le Segreterie operano secondo un orario definito dall'Ateneo che garantisca un minimo di giorni in cui siano accessibili agli studenti; è possibile prendere appuntamento e consultare la Segreteria tramite e-mail e telefonicamente.
            \par
                In aggiunta al servizio di Segreteria, entrambi i Centri monitorano la carriera e l'andamento degli studenti sia in entrata che in uscita. In caso di studenti con titolo estero l'ASTU inoltrerà alle Segreterie la documentazione necessaria, aggiornandole sull'iter di verifica dei Titoli di Studio; una volta accertato, il rispettivo Centro procede all'immatricolazione. L'ASTU viene contattato anche per le richieste delle domande di equipollenza e, in particolare per i CdS Tecnologie digitali e comunicazione per le industrie creative e Comunicazione Multimediale e Tecnologie dell'Informazione, il Centro riceve la documentazione necessaria da inoltrare alle Commissioni Didattiche competenti. \newline
                Sia CEPO che CEGO vengono aggiornati annualmente dall'ASTU sugli eventi di orientamento che si svolgeranno, fornisce loro le informazioni necessarie per poter allestire gli spazi ma lasciando ai Centri l'autonomia sulla logistica. I Centri possono anche proporre e ideare degli eventi a livello locale su richiesta dei CdS attivi. A fine di ciascun anno accademico, i Centri creano un report delle attività da inoltrare al Delegato del Rettore in modo da aggiornare il Magnifico Rettore; il CEGO, dopo aver inoltrato il report, lo utilizza per la stipulazione della Scheda SUA. \newline
                Gli studenti delle Lauree Magistrali, iscritti regolarmente, possono partecipare ad un bando di selezione per poter diventare dei Tutor informativi del DIP a cui afferiscono. Entrambi i Centri poi offrono supporto logistico; il CEGO in aggiunta, monitora i tutor nelle attività di orientamento organizzate dai CdS attivi. 
            \par
                Entrambi i Centri offrono la possibilità di seguire dei tirocini sia interni che esterni a tutti gli studenti iscritti a CdS i quali abbiano nel Curriculum l'attività di tirocinio; la sola differenza tra i due è che il CEGO ha un Ufficio Tirocini proprio mentre l'Ufficio Tirocini di riferimento del CEPO ha sede nel Palazzo Toppo Wassermann a Udine, Ufficio referente anche per la Sede Centrale di Ateneo e il Polo distaccato di Gemona.
            \par
                Tra i vari servizi offerti agli studenti, è presente anche il servizio bibliotecario: entrambe le biblioteche presenti nei Centri, operano in collaborazione con l'ABIB e la rete bibliotecaria di Ateneo.
            \par
                Infine, il CEGO ha un sistema di videosorveglianza a capo dell'Ateneo presso il polo di Santa Chiara mentre il sistema del CEPO, nel polo in Via Prasecco, è a capo del Consorzio Universitario. \newline
                CEPO così come CEGO hanno in dotazione personale formato nel primo soccorso e in prevenzione antincendio, personale che segue periodicamente dei corsi di aggiornamento e partecipa alle prove di evacuazione previste dalla normativa vigente. Essendo il CEGO suddiviso in due poli, il personale formato agisce principalmente nel polo di Santa Chiara ma è a disposizione per il polo in Corso Verdi, dato che quella struttura ha il piano delle emergenze affidato alla Diocesi di Gorizia. Anche il personale formato del CEPO partecipa alle prove d'evacuazione che sono organizzate dal Consorzio Universitario. 